\subsection{R-18: two normal families improve the genuine-padding range}

\paragraph{Statement R-18.}
For $m\geq6$ put
\[
 j_m=\min\{\lfloor m/2\rfloor,m-5\},\qquad
 G_m=2j_m(m-5).
\]
Then
\begin{equation}\label{R18:dimension}
 \dim\mathcal P_{m,4}\geq2m^2+2+G_m.
\end{equation}
For any $n>m\geq6$ satisfying
\begin{equation}\label{R18:padding}
 2n^2\leq2m^2+G_m+2,
\end{equation}
R-14 and G-10 give a four-row multiplicity gap in every
coordinate degree $d\geq d_n$. In addition G-11 gives a
gap for the explicit rectangular partition
$\lambda=(nd/4)^4$ at every degree
$d=E_n k$, $k\geq K_n$, with its displayed effective
constants. In particular these conclusions hold eventually
for $n=\lfloor\rho m\rfloor$ and every fixed
$1<\rho<\sqrt{3/2}$.
The bounds of R-16 and R-18 are both valid; one may use the
larger lower bound at each particular size.

\begin{proof}
Keep the root-of-unity cyclic pencil and notation of R-16,
but take $j_m$ disjoint pairs with starting indices
$\mathcal A=\{0,2,\ldots,2(j_m-1)\}$ and use the entries
\[
 p_a=\varepsilon\ell_{a+1}\quad\hbox{at }(a,a+2),
 \qquad q_a=\varepsilon v_{a+2}\quad\hbox{at }(a+1,a).
\]
Each pair of rows still has support in just its three
columns. Therefore the simultaneous factorization in
R-16 remains exact. It gives two identities:
\begin{align*}
 \ell_a C_{r,a}-v_a C_{r,a+1}-p_a C_{r,a+2}
    &=-2p_aq_a v_a R_{a,r},\\
 q_a C_{r,a}+\ell_{a+1} C_{r,a+1}-v_{a+1} C_{r,a+2}
    &= 2p_aq_a\ell_{a+1}R_{a,r}.
\end{align*}
Both follow by substituting the three two-by-two permanent
factors already displayed in R-16; no other entries occur
in that substitution. Thus for $\varepsilon\ne0$ the
differential image contains the polynomial families
\[
 \begin{aligned}
 U_{a,k}(\varepsilon)
    &=\ell_{a+1}v_a v_{a+2}R_{a,a+k}(\varepsilon),\\
 V_{a,k}(\varepsilon)
    &=\ell_{a+1}^2v_{a+2}R_{a,a+k+1}(\varepsilon),
 \end{aligned}
\]
for $3\leq k\leq m-3$. Their values at zero, using the
same unique complementary matchings, are
\[
 \begin{aligned}
 U_{a,k}&=\ell_{a+1}v_a v_{a+2}
          \prod_{i=a+2}^{a+k-1}v_i
          \prod_{i\notin[a,a+k]}\ell_i,\\
 V_{a,k}&=\ell_{a+1}^2v_{a+2}
          \prod_{i=a+2}^{a+k}v_i
          \prod_{i\notin[a,a+k+1]}\ell_i.
 \end{aligned}
\]
They both have bidegree $(m-k,k)$, and their cyclic
translates are obtained by the same operator $\tau$ as
in R-16.

We prove that the two classes
$\pi_t(U_{0,k})$, $\pi_t(V_{0,k})$ are independent
modulo $\pi_t(W_0)$ whenever
\begin{equation}\label{R18:interior}
 3\leq k,t\leq m-3.
\end{equation}
This will give two independent Vandermonde systems.
Write $q=\zeta$, $K=q^k$, and $X=q^t$ just for the
coefficient calculation below. Retain
$A(x,y)=\prod_{i=k}^{m-1}(x-q^iy)$ from R-16, whose
coefficients $A_r$ are nonzero, and put
\[
 H(z,w)=\prod_{i=2}^{k-2}(z-q^iw),\qquad
 B(z,w)=(z-w)(z-qw)H(z,w).
\]
Thus $F_{0,k}=AB$. Define
\[
 \begin{aligned}
 \alpha&=\frac{K-q}{K-1},&
 \beta&=\frac{q-1}{K-1},\\
 \gamma_2&=\frac{q(K-1)}{qK-1},&
 \gamma_1&=\frac{(q-1)(K-q)}{(K-1)(qK-1)},&
 \gamma_0&=\frac{(q-1)^2}{(K-1)(qK-1)}.
 \end{aligned}
\]
All denominators are nonzero in the range
\eqref{R18:interior}.

The $(x,y)$ factors of $U_{0,k}$ and $V_{0,k}$ are
\[
 A_U=(x-qy)\prod_{i=k+1}^{m-1}(x-q^iy),\qquad
 A_V=(x-qy)^2\prod_{i=k+2}^{m-1}(x-q^iy).
\]
Their coefficient ratios are
\begin{equation}\label{R18:Aratios}
 \frac{(A_U)_r}{A_r}=\alpha q^r+\beta,
 \qquad
 \frac{(A_V)_r}{A_r}
    =\gamma_2 q^{2r}+\gamma_1 q^r+\gamma_0.
\end{equation}
Here is an explicit verification using the finite product
formula from R-16. If $u=q^r$, the first ratio is
$((K-q)u+q-1)/(K-1)$. For the second, put
$D=\prod_{i=k+2}^{m-1}(x-q^iy)$.
The same formula gives
\[
 \frac{D_r}{A_r}
  =\frac{(Ku-1)(qKu-1)}{(K-1)(qK-1)},\quad
 \frac{D_{r-1}}{A_r}
  =\frac{(u-1)(Ku-1)}{(K-1)(qK-1)},\quad
 \frac{D_{r-2}}{A_r}
  =\frac{(u-1)(u-q)}{q(K-1)(qK-1)}.
\]
Consequently $(A_V)_r/A_r$, obtained by subtracting
$2qD_{r-1}$ and adding $q^2D_{r-2}$, has numerator
\[
 (Ku-1)(qKu-1)-2q(u-1)(Ku-1)+q(u-1)(u-q).
\]
Its coefficients of $u^2,u,1$ are respectively
$q(K-1)^2$, $(q-1)(K-q)$, and $(q-1)^2$.
This proves \eqref{R18:Aratios}, including endpoint
coefficients interpreted as zero where necessary.

Let $v_b=z-q^bw$ in the next few displays. In character
$t$, after dividing every coordinate by the nonzero factor
$A_{t-s}$, the four vectors
$\pi_t(zF_{0,k})$, $\pi_t(wF_{0,k})$,
$\pi_t(U_{0,k})$, and $\pi_t(V_{0,k})$ have the coefficient
sequences of $H$ times the following four binary cubics:
\begin{equation}\label{R18:cubics}
 zv_0v_1,\quad wv_0v_1,\quad
 \mathcal U_X=\alpha Xv_{-1}v_1^2+\beta v_0v_2v_{k-1},\quad
 \mathcal V_X=\gamma_2X^2v_0^2v_1
               +\gamma_1Xv_1^2v_{k-1}
               +\gamma_0v_2v_{k-1}v_k.
\end{equation}
To check the shifts in this assertion, the $(z,w)$ factor
of $U_{0,k}$, after replacing $w$ by $q^{-1}w$, is
$v_{-1}v_1^2H$. That factor with $w$ replaced back by
$qw$ becomes $v_0v_2v_{k-1}H$.
The $(z,w)$ factor of $V_{0,k}$, after replacing $w$ by
$q^{-2}w$, is $v_0^2v_1H$; its next two shifts are
$v_1^2v_{k-1}H$ and $v_2v_{k-1}v_kH$.
Together with $q^r=q^{t-s}$ and
\eqref{R18:Aratios}, these identities give exactly
\eqref{R18:cubics}.

These four cubics are linearly independent. Modulo the
first two, it suffices to evaluate at $z/w=1,q$, because
the kernel of those two evaluations is precisely
$v_0v_1\operatorname{span}\{z,w\}$.
Direct substitution and simplification give
\begin{equation}\label{R18:evaldet}
 \begin{vmatrix}
 \mathcal U_X(1,1)&\mathcal V_X(1,1)\\
 \mathcal U_X(q,1)&\mathcal V_X(q,1)
 \end{vmatrix}
 =\frac{(q-K)(q^2-K)(q-1)^6(q+1)}
        {q(K-1)(qK-1)}\ne0.
\end{equation}
The cancellation of $X$ in this expression is exact.
Nonvanishing follows from the order $m$ of $q$ and
$3\leq k\leq m-3$; in particular $q\ne\pm1$ and
$K\notin\{1,q,q^2,q^{-1}\}$.

Under \eqref{R18:interior}, the allowed coefficient
indices $s$ contain at least four consecutive values
$b,b+1,b+2,b+3$, with $0\leq b\leq k-3$.
Multiplication by $H$ takes the cubic monomial basis
to four columns whose restriction to those coordinates
has determinant
\[
 \det(H_{b+u-v})_{0\leq u,v\leq3}.
\]
Up to row and column signs this is the dual
Jacobi--Trudi determinant
\[
 s_{(4^b)}(q^2,q^3,\ldots,q^{k-2}).
\]
It is nonzero by the same alternant argument as in
R-16: there are $N=k-3$ variables and the numerator
exponents for $(4^b)$ are distinct in
$[0,N+3]=[0,k]$, strictly shorter than the order $m$.
For $N=0$, necessarily $b=0$ and the determinant is one.
Thus restriction to the four selected coefficients
preserves independence of the cubics
\eqref{R18:cubics}. This proves the asserted independence
of the two normal classes in every character
\eqref{R18:interior}.

Now suppose, for fixed $k$, that a linear combination
\[
 \sum_{a\in\mathcal A}
        (c_aU_{a,k}+d_aV_{a,k})
\]
belongs to $W_0$. Projecting to character $t$ gives a
linear combination of two independent normal classes,
so both $\sum_a c_aq^{at}$ and
$\sum_a d_aq^{at}$ must vanish.
Since $j_m\leq m-5$, the $j_m$ consecutive values
$t=3,\ldots,j_m+2$ are available. Each of these two
systems is an invertible Vandermonde system in the
distinct $q^a$. Hence every coefficient vanishes.
The $2j_m$ normal vectors are independent for this
bidegree. Different $k$ have different bidegrees, so
the full collection has $2j_m(m-5)=G_m$ independent
classes modulo $W_0$.

Adjoin their polynomial deformations to a basis of
the $2m^2+2$ cyclic Jacobian columns. At zero the
augmented columns are independent. A nonzero coefficient
minor remains nonzero for generic $\varepsilon$, when
all these columns lie in the actual deformed Jacobian
image by the two exact cofactor identities. This proves
\eqref{R18:dimension} by differential rank.

The padding conclusion follows from R-14 and G-10.
For the rectangular conclusion, note that
$G_m\leq m(m-5)$, so \eqref{R18:padding} implies
$n^2<3m^2/2$ and hence $n<4m/3$.
Both hypotheses of G-11 therefore hold.
Finally $G_m=m^2+O(m)$, so the upper admissible ratio
in \eqref{R18:padding} tends to $\sqrt{3/2}$.
\end{proof}

